\documentclass[a4paper,12pt]{article}
\usepackage{../common/macros}
\usepackage{layout}
\usepackage{url}

\title{Webサービス論レポート}
\author{S2620154　吉田 将衛}
\renewcommand{\today}{\the\year 年\the\month 月\the\day 日}
\date{\today}

\begin{document}

\maketitle

% ===== サービス1 =====
\section{Discord}

\subsection{サービス概要}
\begin{itemize}
  \item \textbf{名称：} Discord
  \item \textbf{トップページURL：} \url{https://discord.com}
  \item \textbf{概要：} ゲーマー向けのトークアプリである．サーバをたて，テキストチャットやボイスチャットを，PCやスマートフォンから利用できる．最近では，ゲーム以外のコミュニティも多く利用している．
\end{itemize}

\subsection{選んだ理由}
Discordは，オンラインコミュニケーションのためのプラットフォームであり，特に，ゲーマーの間で利用されており，低レイテンシ，リアルタイム性が求められるサービスであるため，どのような技術的工夫がなされているかに興味を持ったため，選んだ．

\subsection{定量情報}
Discordの会社情報\cite{discord-company}によると，MAUは2億人以上，ゲームをプレイしているユーザは90\%以上，1時間以内にゲームをスタートさせたユーザは最大40\%である．
現在では，2億人以上の月間アクティブユーザが数千ものタイトルで毎月20億時間ものゲームをプレイしている\cite{discord-ceo-2025}．

\subsection{サービスの構成}

Discordのインフラ全体としては，Google Cloud Platformを利用している\cite{discord-cloudflare}．Discordは2015年の創業当時からElixirを使用しており，現在もインフラのバックボーンとして機能している．毎秒2,600万件以上のメッセージを処理している\cite{discord-elixir}．リアルタイム通信において，Erlang VMは最高のツールであるため，Elixirを選択している．しかし，全体の構成としてはハイブリッドである．APIはPythonで書かれており，リアルタイム通信はElixir，メディアストリーミングがC++で書かれており，一部のデータサービスはRustで書かれている．

一方で，ターゲット層がゲーマーであるということも相まって，嫌がらせ目的でのDDoS攻撃を受けることが多い．また，従来と同様の負荷分散では，年間数十万ドルのコストがかかる見込みであり，運用の負担も限界に達していた．そこで，Claudeflare CDNを導入することで，世界中のデータセンターでコンテンツをキャッシュし，さらに，DDoS攻撃自動的に検知し，攻撃トラフィックをブロックすることで，コスト削減とセキュリティの向上を実現した\cite{discord-cloudflare}．
これによって，わずか1年で240万人（前年比9000\%）への急成長を遂げた．また，毎秒200万もの攻撃をCloudflareがブロックすることで，攻撃トラフィックを減らし，ユーザへの影響を最小限に抑えることができた\cite{discord-cloudflare}．

従来のElasticsearchベースのシステムでは，大量のメッセージ検索において，高負荷がかかってしまうという問題があった．そこで，Kubernetesを使用し，管理しやすいサイズの小さなクラスタ群に分割することで，1つのノードが故障することによる影響範囲を最小限に抑えることができるようになった\cite{discord-indexes-trillions}．
20億件を超えるような巨大なコミュニティが出現した．このような巨大コミュニティは従来のようにサーバごとにシャードするだけでは，1つのシャードに収まらなくなる．したがって，巨大サーバを管理するため，専用のマルチシャード構成を持つインデックスを巨大サーバに対して割り当てた．Luceneの最大インデックス容量の制限を回避し，並列処理による検索速度の向上を実現した．

元々Apache Cassandraという分散データベースを使用していた．一般的に，ノード数は数十程度であることがほとんどであるが，データ量の増加に伴って最大クラスターが177ノードまで膨れ上がり，パフォーマンス低下や運用コストの増大に直面した．また，Javaで実装されているため，ガベージコレクションによる一時停止が起こり，レイテンシが増加することもあった．そこで，CassandraをScyllaDBに置き換えることにした．ScyllaDBは，Cassandraと互換性のある分散データベースであり，C++で実装されているため，ガベージコレクションによる一時停止がなく，より高いパフォーマンスを提供できる\cite{discord-infoq-scylladb}．さらに，ScyllaDBへの移行段階で，Rustで書かれたデータサービス層を導入することで，クライアントからのリクエストを集約し，ScyllaDBへのアクセスを効率化できた．結果として，177ノードから72ノードに削減し，突発的な高トラフィックが起こっても，安定したパフォーマンスを維持できるようになった．

Discordでは，250万人の同時音声接続を支えるためにWebRTCを独自にカスタマイズして使用している．基本的には，クライアント同士が直接通信するP2P方式ではなく，サーバを介して通信する方式を採用している．これによりユーザのIPアドレスを隠すことができるため，セキュリティとプライバシーの向上につながっている\cite{discord-voice-webrtc}．
クライアント側での工夫としては，無音声時に帯域はCPUの使用率を削減している．
また，バックエンド側では，音声サーバがクラッシュしたり，DDoS攻撃を受けたりした場合にシステムは自動検知し，別の健全なサーバへ瞬時にユーザを移動させることで，耐障害性を実現している．さらに，世界一3リージョン，30以上のデータセンターに850台以上の音声サーバを配置することで，250万人以上の同時接続，220Gbpsを超えるトラフィックを処理できるようにしている．

% ===== サービス2 =====
\section{GitHub}

\subsection{サービス概要}
\begin{itemize}
  \item \textbf{名称：} GitHub
  \item \textbf{トップページURL：} \url{https://github.com}
  \item \textbf{概要：} ソフトウェア開発のためのプラットフォームである．GitHub上にソースコードをホスティングすることで，数百万人もの他の開発者と一緒にコードレビューやプロジェクト管理を行うことができる．
\end{itemize}

\subsection{選んだ理由}
GitHubは，ソフトウェア開発のためのプラットフォームとして，デファクトスタンダードであり，自分自身も日常的に利用している．また，ソフトウェア工学の研究として，普段からGitHub上にホスティングされたプロジェクトのソースコードを解析することがあるので，どのような技術的工夫がなされているかに興味を持ったため，選んだ．

\subsection{定量情報}
GitHubのブログ\cite{github-octoverse-2025}によると，2025年の時点で，GitHub上には1億8000万人以上の開発者が作業を行なっており，平均して，毎秒1人以上の新規開発者がGitHubに参加している．
開発者は毎分230以上の新しいリポジトリを作成し，毎月平均4320万件のPRをマージ，2025年には10億件のコミットをプッシュした．

\subsection{サービスの構成}
github.com，GitHub Pages，raw.github.comの静的アセットをFastlyのCDNで配信している\cite{github-fastly}．Fastly導入前はリクエストが急増すると，オリジンサーバが過負荷状態になることがあったが，Fastlyの導入によって個々のPagesのWebサイトの表示速度が大幅に向上し，オリジンサーバへのリクエストが最小限に抑えられた．また，比較によると，Arvixe，Windows Azure，Amazon AWSを凌駕するサイトの読み込み速度を実現している\cite{github-fastly}．

GitHubの成長に伴い，データベースの拡張において，様々な問題が起こった．したがって，仮想パーティションという概念を導入した．なぜなら将来的にデータベースを物理的に切り分ける際に問題が発生する可能性があるからである．仮想パーティションは，データベースの論理的な分割を提供し，物理的な切り分けを行う前に，データベースの拡張を可能にする\cite{github-spokes-stretching}．

仮想的に分離されたスキーマドメインを，物理的に別のデータクラスタへ移動するために，Vitessを使用した．Vitessの垂直シャ＝ディング機能を利用して，ダウンタイムなしで複数のテーブルをまとめ，本番環境に移行した．さらなる分割のために水平パーティショニングも使用されている．これによりデータベーステーブルを複数のクラスタに分割できるため，より持続可能な成長が可能になる．

Spokesは3,800万以上のGitリポジトリと3,600万を超えるgistを保存しているファイルサーバのレプリケーションシステムのことを指す．すべてのリポジトリとすべてのgistのコピーを少なくとも3つ保持することで，サーバやネットワークに障害が発生した場合でも永続的で可用性の高いアクセスを提供できる\cite{github-spokes-resilience}．

GitHubは自社のソースコードをgithub.com上でホスティングしている．これにより，もしgithub.comがダウンした場合，それを修正するためのコードにアクセスできず，デプロイもできないという循環依存が発生する．したがって，GitHubはミラーサーバやバックアップを用意しているが，デプロイ用スクリプトの実行時に予期せぬ依存を持つリスクがあった．たとえば，スクリプトの実行中にGitHubからバイナリをダウンロードする動作，ローカルツールが実行時にアップデート確認のためにGitHubへ通信する，あるいはスクリプトが呼び出した別のサービスがGitHubへ通信するといった，直接的，間接的，過渡的な依存が存在することがあった．
このような依存を解決するためのカーネルレベルでフィルタリングする解決策を取った．たとえば，ネットワークの通信を監視したり，DNSクエリを監視し，特定のドメイン，たとえば，github.comへの通信をブロックすることで，スクリプトの実行中にGitHubへの依存がないことを保証できる．
このようにして，GitHubは，github.comへの依存を完全に排除でき，github.comがダウンしても，デプロイ用スクリプトを実行できるようになった．

% ===== 参考文献 =====
\bibliographystyle{jplain}
\bibliography{refs}

% ===== AIツールの使用について（使用した場合のみ残す） =====
\subsection*{AIツールの使用について}
\todo{使用したAIツールの名称と用途を簡潔に記述する．また，本学の「学修におけるChatGPT等の生成AIの利用について」を読んで課題を実施した旨を明記する．使用しなかった場合はこの節を削除する．}

\end{document}
